\section{Method}
\label{sec:method}

\subsection{Problem Formulation}
\label{sec:method-problem}

Given a candidate degrader tuple $\tuple$, where the POI may carry a disease-relevant PTM of type $\tau \in \{\text{phospho}, \text{Ub}, \text{methyl}, \text{acetyl}, \dots\}$ at site $s$, we want to rank candidates by their expected PTM-isoform-selectivity advantage over the wild-type POI. Let $\score(\cdot)$ denote any black-box ternary-complex predictor that returns a real-valued pTernary score for a given $\tuple$ structural instantiation. The raw PTM-selectivity gap is
\begin{equation}
\label{eq:dpt-raw}
\dptraw(\poiptm,\eligase,\protac) \;=\; \score(\poiptm,\eligase,\protac) \;-\; \score(\poiwt,\eligase,\protac).
\end{equation}
The score gap is preferred over the absolute score $\score(\poiptm,\eligase,\protac)$ for two reasons: (i)~PTM-blind scorers are trained on canonical POIs, so absolute scores carry per-POI biases that exceed the PTM signal magnitude, and (ii)~the gap cancels these biases by construction, leaving the change attributable to the PTM perturbation. We will see in \S\ref{sec:experiments-ablation} that omitting this design choice (ranking on $|\score(\poiptm,\cdot)|$) destroys the headline AUC lift; the gap construction is the first load-bearing decision.

\subsection{Per-POI Noise-Floor Calibration}
\label{sec:method-calibration}

Equation~\ref{eq:dpt-raw} on its own is insufficient. Two POIs at different positions in the scorer's training distribution have different intrinsic variance under \emph{any} small perturbation of the POI surface, not just under real PTMs. A raw $\dptraw$ of $0.10$ on one POI may be inside the scorer's noise floor while $0.05$ on a different POI may sit far outside it. Comparing $\dptraw$ across POIs without per-POI normalization conflates real PTM signal with per-POI scorer artifacts.

We resolve this by characterizing the per-POI null distribution before any held-out evaluation. We call this pre-screening pass \textbf{Phase-0} (in analogy to a clinical-trial Phase-0 micro-dose study: a small, fast, low-cost pre-screening run that gates whether the full evaluation budget should be committed). Phase-0 is run once per (scorer, PTM-type) pair on the training subset and its outputs are consumed by every downstream stage. The protocol is:
\begin{enumerate}[leftmargin=*,itemsep=2pt,topsep=2pt]
\item For each POI in the training subset $\trainset$, sample $N{=}200$ random surface positions $\{p_1,\dots,p_{200}\}$ outside the canonical recruitment interface.
\item At each $p_i$, attach a chemically inert mock group whose mass and radius are matched to the target PTM type. The phospho calibration uses $\approx 80$~Da, $\approx 3$~\AA{} radius, neutral, no H-bond donors/acceptors. For mono-ubiquitylation, the mass is $\approx 8.5$~kDa; for methylation, $14$--$42$~Da; for acetylation, $\approx 42$~Da. Mass-matched perturbations protect against false discovery at scorer training-distribution edges.
\item Construct the perturbed POI via constrained side-chain remodel plus a $100$-step minimization; score with the same $\score(\cdot)$ used downstream.
\item Record the per-POI null distribution $\{\dptraw_i = \score(\mathrm{POI}^{(p_i)},\eligase,\protac) - \score(\poiwt,\eligase,\protac)\}_{i=1}^{200}$; fit its mean $\nfmu$ and central 95\% interval $[\nflo,\nfhi]$.
\end{enumerate}

The \emph{calibrated} $\dpt$ score is the raw gap expressed in per-POI noise-floor units, with a clearance threshold:
\begin{equation}
\label{eq:dpt-cal}
\dpt(\poiptm,\eligase,\protac) \;=\;
\begin{cases}
\dfrac{\dptraw(\poiptm,\eligase,\protac)}{|\nfmu|} & \text{if } \big|\dptraw(\poiptm,\eligase,\protac)\big| \;>\; \nftau \cdot |\nfmu|, \\[1.5ex]
0 & \text{otherwise.}
\end{cases}
\end{equation}
The threshold multiplier is pre-registered at $\nftau {=} 1.5$, chosen so that under an approximately-Normal null distribution the clearance criterion accepts only $|\dptraw|$ values that fall in the upper $\sim 5\%$ tail of the per-POI null---i.e., a one-sided $p < 0.05$ rejection of the null hypothesis ``the observed gap is indistinguishable from random surface noise'' for that POI. Tuples below the threshold are reported as ``insufficient signal''; tuples above are ranked by $\dpt$. The cross-track experiments in \S\ref{sec:experiments-scope} use the same $\nftau {=} 1.5$ but with per-PTM-type-matched perturbation mass for the corresponding $\nfmu$ table. \textbf{The methylation track failure (\S\ref{sec:experiments-scope}) is mechanistically traceable to this step}: $14$--$42$~Da methyl perturbations elicit per-POI noise-floor magnitudes that exceed the typical PTM-driven gap, so most methylation-degron tuples are routed to the zero branch of Eq.~\ref{eq:dpt-cal}.

\subsection{Cascaded Pipeline}
\label{sec:method-pipeline}

The full pipeline cascades four black-box modules. Figure~\ref{fig:pipeline} illustrates the data flow.

\begin{figure}[t]
\centering
\begin{tikzpicture}[
  node distance=8mm and 9mm,
  every node/.style={font=\small},
  box/.style={draw, rounded corners=2pt, minimum height=11mm, minimum width=22mm, align=center, fill=blue!4},
  scorer/.style={draw, rounded corners=2pt, minimum height=11mm, minimum width=22mm, align=center, fill=orange!8},
  calib/.style={draw, rounded corners=2pt, minimum height=11mm, minimum width=22mm, align=center, fill=red!6},
  arrow/.style={-Stealth, thick},
  dashed-arrow/.style={-Stealth, thick, dashed, gray}
]
\node[box] (input) {(POI, $\eligase$, \\ $\protac$, $\tau$, $s$)};
\node[box, right=of input] (wt) {Stage A:\\$\poiwt$\\ \scriptsize(PDB / AF-DB v4)};
\node[box, right=of wt] (ptm) {Stage B:\\$\poiptm$ \\ \scriptsize Boltz-2 CCD-PTM};
\node[scorer, right=of ptm] (score) {Stage C:\\ $\score(\cdot)$\\ \scriptsize DeepTernary /\\PROTAC-STAN};
\node[calib, right=of score] (cal) {Stage D:\\ $\dpt$\\ \scriptsize Eq.~\ref{eq:dpt-cal}};
\node[box, below=of ptm, fill=blue!4!yellow!10] (md) {Stage B$'$ (fallback):\\ MD-relaxed\\ \scriptsize ff14SB + phosaa14SB, 50~ns};
\node[calib, below=of cal, fill=red!4!yellow!10] (nf) {Phase-0 table:\\ $\{\nfmu, \nflo, \nfhi\}_{\text{POI}}$\\ \scriptsize Eq.~\ref{eq:dpt-cal} per-POI};

\draw[arrow] (input) -- (wt);
\draw[arrow] (wt) -- (ptm);
\draw[arrow] (ptm) -- (score);
\draw[arrow] (score) -- (cal);
\draw[dashed-arrow] (ptm) -- (md);
\draw[dashed-arrow] (md.east) -- ++(0.5,0) |- (score.west);
\draw[dashed-arrow] (nf) -- (cal);
\node[below=2mm of input, font=\scriptsize, gray] {Input};
\node[below=2mm of cal, font=\scriptsize, gray] {Calibrated ranking};
\end{tikzpicture}
\caption{Pipeline overview. Stages~A and~B produce the wild-type and PTM-occupied POI structures; Stage~C scores both with a black-box ternary predictor; Stage~D computes the noise-floor-calibrated $\dpt$ via Eq.~\ref{eq:dpt-cal}. The MD-relaxed fallback route (Stage~B$'$, dashed) replaces Boltz-2 when CCD-PTM token coverage is missing. The Phase-0 noise-floor table is built per-POI in advance and consumed at Stage~D.}
\label{fig:pipeline}
\end{figure}

\textbf{Stage A (WT structure).} If a PDB or AlphaFold-DB v4 structure for the POI exists with $\mathrm{pLDDT} > 70$ at the PTM site neighborhood, we use it directly. Otherwise, Boltz-2 predicts $\poiwt$ from the canonical UniProt sequence.

\textbf{Stage B (PTM-occupied structure, primary route).} Boltz-2 with Jan-2026 weights ingests the canonical sequence plus a native CCD-PTM token at the experimentally annotated site $s$. We use default sampling settings without Boltz-sample modifications and take the top-ranked seed.

\textbf{Stage B$'$ (PTM-occupied structure, MD fallback).} When CCD-PTM token coverage is missing (rare modifications outside the AF3/Boltz-2 token vocabulary), we start from the Stage~A WT structure, chemically attach the PTM at $t{=}0$, minimize, equilibrate for $1$~ns NPT, then run $50$~ns NPT production with AMBER ff14SB plus phosaa14SB in explicit solvent. The last frame is the PTM-occupied input. We test in \S\ref{sec:experiments-ablation} whether the two routes are comparable within the calibrated $\dpt$ noise; that comparability ($r{\geq}0.7$ between per-tuple route scores) is a deployability premise rather than a strict-equality claim---when both routes are available, the Boltz-2 route remains the head-to-head default.

\textbf{Stage C (scoring).} Both $\poiwt$ and $\poiptm$ are passed to the same scorer $\score(\cdot)$. We instantiate the pipeline with $\score \in \{\text{DeepTernary}, \text{PROTAC-STAN}\}$ to show in \S\ref{sec:experiments-ablation} that the headline $\auclift$ is not specific to a single scorer.

\textbf{Stage D (calibration).} Eq.~\ref{eq:dpt-cal} applies, using the per-POI noise-floor table built once at Phase~0 on $\trainset$.

\subsection{Scorer-Agnostic Instantiation}
\label{sec:method-scorers}

The same Phase-0 protocol is run once per scorer to populate a scorer-specific noise-floor table; downstream tuples are then routed to the table for the scorer in use. For the headline experiment we use DeepTernary; the cross-scorer ablation (\S\ref{sec:experiments-ablation}) substitutes PROTAC-STAN with its own Phase-0 table. Equation~\ref{eq:dpt-cal} is unchanged.

\subsection{Recap of Design Decisions}
\label{sec:method-recap}

The four load-bearing design decisions are: (1)~score gap rather than absolute score (Eq.~\ref{eq:dpt-raw}); (2)~per-POI noise-floor calibration with mass-matched perturbations (Eq.~\ref{eq:dpt-cal}); (3)~clearance threshold $\nftau{=}1.5$ pre-registered before held-out evaluation; (4)~MD-relaxed fallback to decouple from PTM-conditioned ensemble prediction. The first three are evaluated as ablations in \S\ref{sec:experiments-ablation}; the fourth is evaluated as a route-equivalence check.
